\chapter{Solid State Theory: Undated Midterm Fragment}
\label{q:theory-undated-midterm-fragment}
\begin{sourceNote}
Two photographs preserve the end of Question 2 and Questions 3--7. The date, Question 1, and the model and parameters at the start of Question 2 are outside the photographed area. The visible weights total 85\%. No missing wording has been reconstructed from guesses or handwritten answers. The source calls the single-layer honeycomb model ``two-dimensional graphite''; this terminology is retained below.
\end{sourceNote}

\section*{2. Band gaps and effective mass: incomplete question}
\emph{Only the following instruction and subparts are visible:} Estimate the following quantities:
\begin{enumerate}[label=(\alph*)]
\item (10\%) The band gap between the lowest two bands.
\item (5\%) The band gap between the second and the third lowest bands.
\item (5\%) The effective mass for the states about the top of the lowest band.
\end{enumerate}
\emph{The defining potential, model, and numerical parameters are missing.}

\section*{3. Tight-binding dispersion (10\%)}
Derive the band dispersion
\[
E_k=E_{k=0}+2\gamma(\cos ka-1)
\]
according to the tight-binding model, where $a$ is the lattice constant and $\gamma<0$ is the inter-orbital overlap (the hopping parameter).

\section*{4. True or false (10\%)}
Explain your answer.
\begin{enumerate}[label=(\alph*)]
\item (5\%) Two atoms of the same type in a crystal must be connected by a translation vector of the underlying Bravais lattice.
\item (5\%) Consider a $\vct k$ point lying outside the first Brillouin zone of a Bravais lattice. There must exist a reciprocal lattice point to which the $\vct k$ point is closer than to the origin.
\end{enumerate}

\section*{5. Lattice planes (15\%)}
Consider a Bravais lattice spanned by the primitive vectors
\[
\vct t_1=\frac a2(0,1,1),\qquad
\vct t_2=\frac a2(1,0,1),\qquad
\vct t_3=\frac a2(1,1,0).
\]
Here $\vct g_i$ denote the reciprocal primitive vectors.
\begin{enumerate}[label=(\alph*)]
\item (5\%) Consider a family of direct lattice planes normal to $\tfrac12\vct g_1+\vct g_2-\vct g_3$. Find the distance between two successive lattice planes.
\item (5\%) The lattice points on each plane in (a) form a two-dimensional lattice. Find the number of lattice points per unit area on that plane.
\item (5\%) Consider a family of direct lattice planes normal to $\vct t_1+2\vct t_2-2\vct t_3$. Find the distance between two successive lattice planes.
\end{enumerate}

\section*{6. Two-dimensional graphite (25\%)}
Consider the two-dimensional crystal of graphite, which has two carbon atoms per unit cell at positions
\[
\vct d_1=0,\qquad \vct d_2=a(0,1/\sqrt3).
\]
The primitive translation vectors for the underlying two-dimensional hexagonal lattice are
\[
\vct t_1=a(1/2,\sqrt3/2),\qquad
\vct t_2=a(-1/2,\sqrt3/2).
\]
Graphite has two $\pi$ bands with dispersion
\[
\varepsilon_{\vct k}=\varepsilon_p\pm\gamma|F_{\vct k}|,
\qquad
|F_{\vct k}|^2=1+4\cos\frac{k_xa}{2}\cos\frac{\sqrt3k_ya}{2}
+4\cos^2\frac{k_xa}{2},
\]
where $\varepsilon_p$ is constant and $\gamma>0$.
\begin{enumerate}[label=(\alph*)]
\item (15\%) Identify whether the symmetry points
\[
\Gamma:\ \vct k=0,\qquad
Q:\ \vct k=\frac{2\pi}{a}\left(\frac12,\frac1{2\sqrt3}\right),\qquad
P:\ \vct k=\frac{2\pi}{a}\left(\frac23,0\right)
\]
of the first Brillouin zone are critical points. What types of critical points are they?
\item (10\%) Express the density of states contributed by states in a small region containing a critical point of your choice from (a).
\end{enumerate}

\section*{7. Bloch theorem (5\%)}
Describe the Bloch theorem.
